%% find-evil / VERITAS — Project Story (Devpost narrative)
%% SANS DFIR Hackathon 2026 | Category 7: Persistent Learning Loop

\documentclass[10pt, letterpaper]{article}

\usepackage[margin=0.75in, top=0.65in, bottom=0.65in]{geometry}
\usepackage{titlesec}
\usepackage{parskip}
\usepackage{microtype}
\usepackage{xcolor}
\usepackage{fontenc}
\usepackage{mdframed}
\usepackage{siunitx}
\usepackage{hyperref}
\hypersetup{colorlinks=true, urlcolor=blue!60!black, linkcolor=black}

%% ── Colour palette ─────────────────────────────────────────────────────────
\definecolor{accent}{RGB}{29, 78, 216}
\definecolor{redaccent}{RGB}{220, 38, 38}
\definecolor{codebg}{RGB}{248, 250, 252}

%% ── Section style ───────────────────────────────────────────────────────────
\titleformat{\section}[block]
  {\normalfont\bfseries\color{accent}}
  {}{0pt}{}[{\color{accent}\titlerule[0.8pt]}]
\titlespacing{\section}{0pt}{8pt}{4pt}

%% ── Stat box ────────────────────────────────────────────────────────────────
\newmdenv[
  backgroundcolor=codebg,
  linecolor=accent,
  linewidth=1pt,
  leftmargin=0pt, rightmargin=0pt,
  innerleftmargin=8pt, innerrightmargin=8pt,
  innertopmargin=5pt, innerbottommargin=5pt,
  skipabove=6pt, skipbelow=6pt
]{statbox}

%%─────────────────────────────────────────────────────────────────────────────
\begin{document}

%% ── Header ──────────────────────────────────────────────────────────────────
\begin{center}
  {\LARGE\bfseries\color{accent} find-evil / VERITAS}\\[3pt]
  {\large Adversarial Forensic Verification for Windows Incident Response}\\[2pt]
  {\small SANS DFIR Hackathon 2026 \quad|\quad Category 7: Persistent Learning Loop}
\end{center}

\begin{statbox}
\small\centering
\textbf{4} APT machines confirmed compromised\quad
\textbf{2} triage false positives caught and refuted on physical evidence\quad
\textbf{100\%} of confirmed findings verified against disk artifacts\quad
\textbf{4-layer} MCP security boundary\quad
\textbf{9} MITRE techniques covered by corpus-calibrated signal weights
\end{statbox}

%% ── Sections ────────────────────────────────────────────────────────────────
\section{Inspiration}

Every LLM-based forensic tool has the same problem: the model
\emph{wants} to find evidence. It is helpful by design. Give it a
disk image and ask whether credential dumping occurred, and it will
find something that looks like credential dumping---whether or not the
binary is actually on disk.

The standard answer is prompt engineering: tell the model to be
careful, to be skeptical, to only report confirmed findings. Prompt
controls are not security controls. They can be overridden, forgotten,
or simply ignored when the model is confident.

VERITAS answers a different question: what if the architecture
\emph{itself} made hallucinating a confirmed finding structurally
impossible? Not unlikely. Impossible. A finding is only confirmed when a
second independent agent---one instructed to distrust the first---calls
a forensic tool and reads the actual bytes off the disk. If the file is
not there, the technique is refuted. No amount of model confidence
changes that.

That is the adversarial dynamic in VERITAS: not Red vs.\ Blue in
training, but Investigator vs.\ Auditor in investigation. The Disk Agent
and Memory Agent investigate. The Forensic Auditor demands proof.

\section{What It Does}

VERITAS investigates any mounted Windows forensic image through a
three-phase pipeline, fully autonomous from invocation to HTML report.

\textbf{Phase 1 — Deterministic triage.} The Disk Agent dispatches
approximately 25 generic SIFT commands in under 60 seconds---no LLM,
no case-specific assumptions---and scores the image against
corpus-calibrated signal weights across 9 MITRE ATT\&CK techniques.
Weights are log-odds ratios derived from 800+ labeled malware samples.
Every command is invariant across investigations: nothing from a previous
case contaminates the baseline sweep.

\textbf{Phase 2 — Agentic deep investigation.} A Claude-powered
investigation loop with a 25-call tool budget follows up on triage
findings. The agent receives an explicit list of what was already
checked and a directed list of uncovered investigation domains---event
log content, prefetch binary parsing, shellbags, SAM/SECURITY hive
extraction, hash verification. It cannot re-run what Pass~1 already
covered; every call advances the investigation into new territory.

\textbf{Phase 3 — Forensic Auditor.} After investigation completes,
the Forensic Auditor challenges every detected technique in parallel
(\texttt{asyncio.gather}), running up to five rounds of two tool calls
each. The rule is simple: identify the artifact that \emph{must}
physically exist on disk if this technique executed, then look for it.
Signal string-match alone never confirms a finding. A finding is
\textsc{Confirmed} only when the definitive artifact is present;
\textsc{Refuted} when the Auditor finds positive evidence of absence.

Multi-host campaigns are supported natively. Confirmed IOCs from one
investigation feed the next via \texttt{--ioc-file}, and
\texttt{adversa-merge-iocs.sh} merges findings across hosts into a
unified campaign IOC set for forward-deployment on new images.

\section{How We Built It}

\textbf{One tool, four security layers.} Every forensic action flows
through a single MCP primitive: \texttt{run\_terminal\_command}. Behind it
is a four-layer validator. 22 hard-blocked dangerous strings (\texttt{sudo},
command substitution, \texttt{shred}, network tools). A 53-tool SIFT
binary allowlist. A quote-aware pipeline parser---necessary because
standard forensic invocations like \texttt{grep -iE '(http|https|ftp)'}
contain \texttt{|} inside the pattern argument, and a naive split would
reject \texttt{https} as an unlisted binary. A redirect guard that
verifies all output lands in \texttt{reports/} and nowhere else.
Evidence modification is structurally impossible, not just unlikely.

\textbf{Append-only audit log.} Every tool call is written to the audit
log \emph{before} it executes: the command, the agent that called it,
a timestamp, and whether the validator passed or blocked it. Chain of
custody is not a report we produce afterward---it is a property of the
system by construction.

\textbf{Plain-text verifiability.} All Auditor output is plain prose.
Every confirmed finding in the HTML report cites the exact tool call
that produced it. A reviewer can open \texttt{reports/audit\_log.jsonl},
find the entry, and reproduce the result with one shell command on the
same mounted image.

\textbf{Corpus-calibrated signal weights.} Detection signals are weighted
using log-odds ratios computed from 800+ labeled malware samples sourced
from MalwareBazaar and HybridAnalysis. Each signal's weight reflects
how discriminative it is between malware and a benign Windows baseline.
Cross-technique tokens are dampened (IDF-equivalent); signals from confirmed
cases retain a floor weight. No neural network---every weight is a number
with a traceable source sample.

\section{Challenges We Ran Into}

The hardest engineering problem was the validator rejecting legitimate
forensic commands. The first version used a regex to split on \texttt{|}
and checked each segment's leading binary against the allowlist. The
first time the Disk Agent ran \texttt{grep -iE '(http|https|ftp)'},
the validator split it into three segments and rejected \texttt{https}
as an unlisted binary. Fixing this required a quote-aware parser that
tracks single-quoted substrings and correctly identifies \texttt{|} inside
a quoted argument as part of the argument, not a pipeline separator.

The harder intellectual problem was false positives. Pass~1 triage on the
controller returned a score of 145 with three detected techniques.
Two of them were wrong: masquerading signals triggered by legitimate
\texttt{svchost.exe} copies in WinSxS, and account discovery signals
triggered by a user profile \emph{directory} rather than active
enumeration tools. Without the Auditor, an analyst would have received
three work items. With the Auditor, they received one---the only technique
that left a physical artifact.

\section{Accomplishments That We're Proud Of}

Four real APT machines from a single intrusion, investigated
autonomously with a consistent pipeline and a complete audit trail.

nfury: T1003.001 confirmed with \texttt{hydrakatz.exe} found on disk and
hash-verified; T1087.001 confirmed via SAM hive analysis. Score 95, no
false positives.

nfury: 19 techniques investigated, 15 confirmed, 4 refuted by the Auditor. Four
false positives refuted on physical evidence---the exact distinction between
``this binary exists in a legitimate Windows location'' and ``this binary
was placed there by an attacker.'' One technique confirmed:
\texttt{procdump.exe} in \texttt{/Tools/SysInternals/} and
\texttt{spinlock.exe} in the WER ReportQueue crash dump, confirming the
implant executed on this machine.

tdungan: Investigated as a blind forward-validation with campaign IOCs
from the merged nfury and controller investigation. T1003.001 and
T1071.001 both confirmed with physical artifact verification. Triage 67,
adjusted to 85 after Auditor IOC cross-match. No false positives.

The \texttt{vibranium} domain account and \texttt{spinlock.exe} implant
were identified across the enterprise from signals that were extracted
from a completely different dataset---Mordor Sysmon telemetry from 2019,
not from these images.

\section{What We Learned}

Physical artifact verification is not a feature. It is the only thing
that separates a forensic finding from a model opinion. Every
architecture decision in VERITAS follows from that single principle.

Architectural guardrails beat prompt controls. The four-layer MCP
boundary means no version of the model can write to evidence directories,
spawn network connections, or execute arbitrary shell code---not because
we told it not to, but because those paths don't exist in the validator.
The Forensic Auditor's refutation of false positives works the same way:
the model cannot confirm a technique without calling a tool and reading
actual output. Confidence is not evidence.

The most valuable component is the one nobody asked for. The Auditor was
not in the original design. Adding a second agent whose only job is to
disprove the first agent's findings turned a triage tool into an
investigation system.

\section{What's Next for VERITAS}

The detection rule foundation needs pool-separation: a signal should only
be admitted if it appears in documented attack telemetry and \emph{never}
in a real benign baseline. This eliminates the false signal problem at
training time rather than relying on the Auditor to catch it at
investigation time.

Beyond that: memory forensics integration (Volatility~3 via the MCP
layer, for process injection and rootkits invisible on disk), timeline
correlation (Plaso super-timeline filtered to confirmed technique time
windows), and coverage expansion from 8 to 50+ MITRE techniques via
automatic Mordor dataset enumeration.

\end{document}
